\documentclass[11pt]{article}

% Official *ACL style, used by ARR and every *ACL venue. Use [review] for an
% anonymous ARR submission and [final] for camera-ready.
\usepackage[preprint]{acl}
\usepackage{times}
\usepackage{latexsym}
\usepackage[T1]{fontenc}
\usepackage[utf8]{inputenc}
\usepackage{microtype}
\usepackage{amsmath,amssymb}
\usepackage{booktabs}
\usepackage{graphicx}

\title{LATTICE_PROJECT_TITLE}
\author{Author Name \\
  Institution \\
  \texttt{author@example.com}}

\begin{document}
\maketitle

\begin{abstract}
Briefly summarize the problem, method, and main result.
\end{abstract}

\section{Introduction}
Motivate the problem and state the paper's main contribution.

\section{Related Work}
Position the contribution against the most relevant prior work.

\section{Method}
Describe the proposed method.

\section{Experiments}
Present the evaluation setup and results.

\section{Conclusion}
Summarize the findings.

\section*{Limitations}
State the limitations of the work. ARR requires this section; it does not
count toward the page limit.

% acl.sty stops the build on an empty reference list, which a new paper has
% until its first \cite. Uncomment this line with the first citation.
% \bibliography{references}

\end{document}
